\documentclass[11pt,UTF8,fontset=fandol]{ctexart}
\usepackage{amsmath,amssymb,amsfonts}
\usepackage{graphicx,float,caption,subcaption}
\usepackage{booktabs,array,multirow,makecell}
\usepackage{geometry,enumitem,xcolor}
\usepackage[hidelinks]{hyperref}
\usepackage{url}
\usepackage{longtable}
\usepackage{listings}
\lstdefinestyle{promptstyle}{
  basicstyle=\ttfamily\footnotesize,
  breaklines=true,
  breakatwhitespace=false,
  columns=fullflexible,
  keepspaces=true,
  frame=single,
  framesep=4pt,
  xleftmargin=4pt,
  xrightmargin=4pt,
  aboveskip=6pt,
  belowskip=6pt,
  breakindent=0pt,
  postbreak=\mbox{},
}
\geometry{margin=2.3cm}
\graphicspath{{/Users/yye/project/论文/assets/graphrag/}}
\captionsetup{font=small,labelfont=bf}
\setlength{\parskip}{0.4em}

\title{\textbf{从局部到全局：一种用于查询聚焦摘要的 GraphRAG 方法}}

\author{Darren Edge$^{1\dagger}$ \quad Ha Trinh$^{1\dagger}$ \quad Newman Cheng$^{2}$ \quad Joshua Bradley$^{2}$ \quad Alex Chao$^{3}$\\[2pt]
Apurva Mody$^{3}$ \quad Steven Truitt$^{2}$ \quad Dasha Metropolitansky$^{1}$ \quad Robert Osazuwa Ness$^{1}$\\[2pt]
Jonathan Larson$^{1}$\\[4pt]
$^{1}$\,Microsoft Research \quad $^{2}$\,Microsoft Strategic Missions and Technologies \quad $^{3}$\,Microsoft Office of the CTO\\[4pt]
\texttt{\{daedge,trinhha,newmancheng,joshbradley,achao,moapurva,}\\
\texttt{steventruitt,dasham,robertness,jolarso\}@microsoft.com}\\[4pt]
{\small $^{\dagger}$\,这些作者对本工作的贡献相同。}}

\date{}
\begin{document}
\maketitle

\begin{abstract}
利用检索增强生成（retrieval-augmented generation, RAG）从外部知识源检索相关信息，能够使大语言模型（large language models, LLMs）回答关于私有和/或此前未见过的文档集合的问题。然而，RAG 在面向整个文本语料库的全局性问题（例如"该数据集的主要主题是什么？"）时会失效，因为这本质上是一个查询聚焦摘要（query-focused summarization, QFS）任务，而非显式的检索任务。与此同时，已有的 QFS 方法又无法扩展到典型 RAG 系统所索引的文本规模。为了结合这两类对比鲜明的方法各自的优势，我们提出 GraphRAG，一种基于图的私有文本语料库问答方法，它能够随用户问题的概括性与源文本的数量一同扩展。我们的方法使用一个大语言模型分两个阶段构建图索引（graph index）：首先，从源文档中导出一个实体知识图谱（entity knowledge graph）；然后，为所有紧密相关实体的分组预先生成社区摘要（community summaries）。给定一个问题，先用每个社区摘要生成一个局部响应，然后将所有局部响应再次汇总为返回给用户的最终响应。对于针对百万 token 量级数据集的一类全局意义建构（global sensemaking）问题，我们的实验表明，相比传统 RAG 基线，GraphRAG 在生成答案的全面性（comprehensiveness）与多样性（diversity）两方面均带来了显著提升。
\end{abstract}

\section{引言}

检索增强生成（RAG）（Lewis et al., 2020）是一种成熟的方法，它利用大语言模型（LLMs）来回答基于数据的查询，而这些数据规模过大，无法完全容纳于语言模型的上下文窗口（context window）之中——上下文窗口指 LLM 一次能够处理的最大 token（文本单位）数量（Kuratov et al., 2024; Liu et al., 2023）。在标准的 RAG 设置中，系统可以访问一个庞大的外部文本记录语料库，并检索出一个记录子集，这些记录单独来看与查询相关，整体上又足够小，能够纳入 LLM 的上下文窗口。随后，LLM 基于查询与所检索到的记录共同生成响应（Baumel et al., 2018; Dang, 2006; Laskar et al., 2020; Yao et al., 2017）。我们将这种传统方法统称为向量 RAG（vector RAG），它在能够利用一小组记录内局部化信息回答的查询上表现良好。然而，向量 RAG 方法并不支持意义建构（sensemaking）查询，即那些需要对整个数据集进行全局理解的查询，例如"在过去十年中，科学发现受跨学科研究影响的关键趋势是什么？"

意义建构任务要求对"连接（这些连接可以存在于人、地点和事件之间）进行推理，以便预判其发展轨迹并采取有效行动"（Klein et al., 2006）。诸如 GPT（Achiam et al., 2023; Brown et al., 2020）、Llama（Touvron et al., 2023）和 Gemini（Anil et al., 2023）之类的大语言模型，在科学发现（Microsoft, 2023）和情报分析（Ranade and Joshi, 2023）等复杂领域的意义建构方面表现出色。给定一个意义建构查询以及一段包含隐式且相互关联的概念集合的文本，LLM 能够生成回答该查询的摘要。然而，挑战在于：当数据量大到需要采用 RAG 方法时，向量 RAG 方法无法支持对整个语料库的意义建构。

本文提出 GraphRAG——一种基于图的 RAG 方法，能够对整个大型文本语料库进行意义建构。GraphRAG 首先使用 LLM 构建一个知识图谱（knowledge graph），其中节点对应语料库中的关键实体，边表示这些实体之间的关系。接着，它将该图划分为紧密相关实体所构成的社区（communities）层级结构，然后使用 LLM 生成社区级别的摘要。这些摘要按照所提取社区的层级结构以自底向上（bottom-up）的方式生成，层级较高的摘要会递归地纳入较低层级的摘要。这些社区摘要共同提供了对语料库的全局描述与洞见。最后，GraphRAG 通过对社区摘要进行 map-reduce 处理来回答查询：在 map 步骤中，各社区摘要被独立且并行地用于为查询提供局部答案；随后在 reduce 步骤中，这些局部答案被合并并用于生成最终的全局答案。

GraphRAG 方法及其对整个语料库进行全局意义建构的能力构成了本工作的主要贡献。为了证明这一能力，我们开发了 LLM-as-a-judge 技术（Zheng et al., 2024）的一种新颖应用，适用于针对广泛议题与主题、且不存在标准答案（ground-truth answer）的问题。该方法首先使用一个 LLM，基于语料库特定的用例生成一组多样化的全局意义建构问题，然后使用第二个 LLM，依据预先定义的标准（在第~3.3~节中定义）来评判两个不同 RAG 系统的答案。我们使用该方法在两个具有代表性的真实世界文本数据集上对 GraphRAG 与向量 RAG 进行比较。结果表明，当使用 GPT-4 作为 LLM 时，GraphRAG 明显优于向量 RAG。

GraphRAG 已作为开源软件发布于 \url{https://github.com/microsoft/graphrag}。此外，GraphRAG 方法的若干版本也已作为扩展集成到多个开源库中，包括 LangChain（LangChain, 2024）、LlamaIndex（LlamaIndex, 2024）、NebulaGraph（NebulaGraph, 2024）和 Neo4J（Neo4J, 2024）。

\section{背景}

\subsection{RAG 方法与系统}

RAG 通常泛指这样一类系统：使用用户查询从外部数据源检索相关信息，随后将该信息纳入由 LLM（或其他生成式 AI 模型，如多媒体模型）对查询所生成响应的过程之中。查询与检索到的记录共同填充一个提示模板（prompt template），随后该模板被传递给 LLM（Ram et al., 2023）。当数据源中的记录总量过大、无法纳入对 LLM 的单次提示中（即数据源中的文本量超过 LLM 的上下文窗口）时，RAG 是理想的选择。在标准的 RAG 方法中，检索过程返回一定数量、在语义上与查询相似的记录，而生成的答案仅使用这些检索到的记录中的信息。传统 RAG 的一种常见方法是使用文本嵌入（text embeddings），在向量空间中检索与查询最接近的记录，其中"接近"对应于语义相似性（Gao et al., 2023）。尽管某些 RAG 方法可能采用其他的检索机制，我们仍将这一类传统方法统称为向量 RAG。GraphRAG 与向量 RAG 的区别在于，它能够回答那些需要对整个数据语料库进行全局意义建构的查询。

GraphRAG 建立在已有的高级 RAG 策略工作之上。GraphRAG 将对数据源大段内容的摘要作为一种"自记忆"（self-memory）形式加以利用（如 Cheng et al. 2024 所述），这些摘要随后被用于回答查询（如 Mao et al. 2020）。这些摘要被并行生成，并迭代式地聚合为全局摘要，这与已有技术类似（Feng et al., 2023; Gao et al., 2023; Khattab et al., 2022; Shao et al., 2023; Su et al., 2020; Trivedi et al., 2022; Wang et al., 2024）。具体而言，GraphRAG 与其他使用层级索引（hierarchical indexing）来创建摘要的方法类似（类似于 Kim et al. 2023; Sarthi et al. 2024）。GraphRAG 与这些方法的区别在于：它从源数据生成一个图索引，然后应用基于图的社区检测（community detection）来对数据进行主题性划分。

\subsection{将知识图谱与 LLM 及 RAG 结合使用}

从自然语言文本语料库中抽取知识图谱的方法包括规则匹配、统计模式识别、聚类和嵌入（Etzioni et al., 2004; Kim et al., 2016; Mooney and Bunescu, 2005; Yates et al., 2007）。GraphRAG 属于较新近的一类研究，这类研究使用 LLM 进行知识图谱抽取（Ban et al., 2023; Melnyk et al., 2022; OpenAI, 2023; Tan et al., 2017; Trajanoska et al., 2023; Yao et al., 2023; Yates et al., 2007; Zhang et al., 2024a）。它也为日益增多的、将知识图谱用作索引的 RAG 方法添砖加瓦（Gao et al., 2023）。一些技术直接在提示中使用子图、图的元素或图结构的属性（Baek et al., 2023; He et al., 2024; Zhang, 2023），或将其作为生成输出的事实依据（Kang et al., 2023; Ranade and Joshi, 2023）。另一些技术（Wang et al., 2023b）使用知识图谱来增强检索，在查询时由基于 LLM 的智能体（agent）动态地遍历一个图，图中节点表示文档元素（例如段落、表格），边则编码词法和语义上的相似性或结构关系。GraphRAG 与这些方法的区别在于，它聚焦于图在此情境下一个此前未被探索的特性：图固有的模块性（modularity）（Newman, 2006），以及将图划分为由紧密相关节点构成的嵌套模块化社区的能力（例如 Louvain，Blondel et al. 2008；Leiden，Traag et al. 2019）。具体而言，GraphRAG 通过使用 LLM 创建跨越该社区层级的摘要，递归地生成越来越全局化的摘要。

\subsection{面向 RAG 评估的自适应基准测试}

目前已有许多面向开放域问答的基准数据集，包括 HotPotQA（Yang et al., 2018）、MultiHop-RAG（Tang and Yang, 2024）和 MT-Bench（Zheng et al., 2024）。然而，这些基准面向的是向量 RAG 的性能，也就是说，它们评估的是在显式事实检索上的表现。在本工作中，我们提出一种方法，用以生成一组问题来评估对整个语料库的全局意义建构能力。我们的方法与那些使用语料库生成问题（其答案应当是对语料库的摘要）的 LLM 方法相关，例如 Xu and Lapata（2021）。然而，为了产生公平的评估，我们的方法避免直接从语料库本身生成问题（作为另一种实现方式，可以使用从后续图抽取和答案评估步骤中保留出来的语料库子集）。

自适应基准测试（adaptive benchmarking）指的是动态生成针对特定领域或用例量身定制的评估基准的过程。近期的工作已使用 LLM 进行自适应基准测试，以确保与目标应用或任务的相关性、多样性和一致性（Yuan et al., 2024; Zhang et al., 2024b）。在本工作中，我们提出一种自适应基准测试方法，用以为 LLM 生成全局意义建构查询。我们的方法建立在已有的基于 LLM 的人物画像（persona）生成工作之上，即使用 LLM 生成多样且真实的人物画像集合（Kosinski, 2024; Salminen et al., 2024; Shin et al., 2024）。我们的自适应基准测试流程使用人物画像生成来创建能够代表真实世界 RAG 系统使用情况的查询。具体而言，我们的方法使用 LLM 推断可能使用该 RAG 系统的潜在用户及其用例，并以此为指导生成语料库特定的意义建构查询。

\subsection{RAG 评估标准}

我们的评估依赖 LLM 来评估 RAG 系统回答所生成问题的效果如何。已有工作表明，LLM 是自然语言生成的良好评估者，其中包括 LLM 评估结果可与人工评估相媲美的工作（Wang et al., 2023a; Zheng et al., 2024）。一些先前的工作提出了让 LLM 量化所生成文本质量的标准，例如"流畅性"（fluency）（Wang et al., 2023a）。这些标准中有一部分是面向向量 RAG 系统的通用标准，与全局意义建构无关，例如"上下文相关性"（context relevance）、"忠实性"（faithfulness）和"答案相关性"（answer relevance）（RAGAS，Es et al. 2023）。在缺乏评估金标准（gold standard）的情况下，可以通过提示 LLM 比较来自两个相互竞争模型的生成结果，来量化给定标准下的相对性能（LLM-as-a-judge，（Zheng et al., 2024））。在本工作中，我们设计了用于评估 RAG 针对全局意义建构问题所生成答案的标准，并使用该比较方法来评估我们的结果。我们还使用从 LLM 抽取的可验证事实陈述（即"声明"，claims）所导出的统计量来验证结果。

\begin{figure}[H]
\centering
\includegraphics[width=0.95\linewidth]{cand_p04_k00.pdf}
\caption{使用源文档文本的 LLM 导出图索引的 GraphRAG 流程。该图索引涵盖节点（例如实体）、边（例如关系）和协变量（covariates，例如声明），它们由针对数据集所属领域定制的 LLM 提示来检测、抽取和摘要。社区检测（例如 Leiden，Traag et al., 2019）被用于将图索引划分为若干元素组（节点、边、协变量），LLM 可以在索引时（indexing time）和查询时（query time）对其进行并行摘要。给定查询的"全局答案"是通过对所有报告了与该查询相关性的社区摘要进行最后一轮查询聚焦摘要而产生的。图中各阶段标注为：\textbf{索引时}（Indexing Time）、\textbf{流程阶段}（Pipeline Stage）、\textbf{查询时}（Query Time）；处理流程包括：源文档（Source Documents）经文本抽取与分块（text extraction and chunking）得到文本块（Text Chunks），经领域定制摘要得到实体与关系（Entities \& Relationships），再得到知识图谱（Knowledge Graph），经社区检测（community detection）得到图社区（Graph Communities），经领域定制摘要得到社区摘要（Community Summaries），经查询聚焦摘要得到社区答案（Community Answers），最终经查询聚焦摘要得到全局答案（Global Answer）。}
\label{fig:1}
\end{figure}

\section{方法}

\subsection{GraphRAG 工作流}

图~\ref{fig:1} 展示了 GraphRAG 方法与流程的高层数据流。在本节中，我们描述每个步骤的关键设计参数、技术和实现细节。

\subsubsection{源文档 → 文本块}

首先，语料库中的文档被切分为文本块（text chunks）。LLM 从每个文本块中抽取信息，以供下游处理。选择文本块的大小是一项基本的设计决策；较长的文本块进行此类抽取所需的 LLM 调用次数更少（从而降低成本），但会因块中靠前位置出现的信息召回率（recall）下降而受到影响（Kuratov et al., 2024; Liu et al., 2023）。有关提示以及召回率—精确率权衡的示例，参见第~A.1~节。

\subsubsection{文本块 → 实体与关系}

在这一步中，LLM 被提示从给定文本块中抽取重要实体的实例以及实体之间的关系。此外，LLM 还会为这些实体和关系生成简短的描述。举例来说，假设某个文本块包含以下文本：

\begin{quote}
NeoChip（NC）的股票在 NewTech Exchange 上市交易的第一周大幅飙升。然而，市场分析师告诫称，这家芯片制造商的公开上市可能并不能反映其他科技 IPO 的趋势。NeoChip 此前是一家私营企业，于 2016 年被 Quantum Systems 收购。这家创新的半导体公司专注于面向可穿戴设备和物联网（IoT）设备的低功耗处理器。
\end{quote}

LLM 被提示抽取出以下内容：

\begin{itemize}[leftmargin=2em]
\item 实体 NeoChip，其描述为"NeoChip 是一家公开上市的公司，专注于面向可穿戴设备和物联网设备的低功耗处理器。"
\item 实体 Quantum Systems，其描述为"Quantum Systems 是一家此前拥有 NeoChip 的公司。"
\item NeoChip 与 Quantum Systems 之间的一种关系，其描述为"Quantum Systems 从 2016 年起拥有 NeoChip，直到 NeoChip 公开上市为止。"
\end{itemize}

通过为上下文学习（in-context learning）（Brown et al., 2020）选择领域适配的少样本范例（few-shot exemplars），这些提示可以针对文档语料库所属的领域进行定制。例如，虽然我们默认的提示抽取的是诸如人物、地点和组织等广泛类别的"命名实体"（named entities）且普遍适用，但具有专业知识的领域（例如科学、医学、法律）将受益于针对这些领域专门设计的少样本范例。

LLM 还可以被提示抽取关于所检测实体的声明（claims）。声明是关于实体的重要事实陈述，例如日期、事件以及与其他实体的交互。与实体和关系一样，上下文学习范例可以提供领域特定的指导。从上述示例文本块中抽取的声明描述如下：

\begin{itemize}[leftmargin=2em]
\item NeoChip 的股票在 NewTech Exchange 上市交易的第一周大幅飙升。
\item NeoChip 作为一家公开上市公司在 NewTech Exchange 首次亮相。
\item Quantum Systems 于 2016 年收购了 NeoChip，并持有其所有权直至 NeoChip 上市。
\end{itemize}

有关实体与声明抽取实现的提示和细节，参见附录~A。

\subsubsection{实体与关系 → 知识图谱}

使用 LLM 抽取实体、关系和声明是抽象式摘要（abstractive summarization）的一种形式——这些是对概念的有意义的摘要，就关系和声明而言，它们可能并未在文本中显式陈述。实体/关系/声明抽取过程会为单个元素创建多个实例，因为某个元素通常会在多篇文档中被多次检测和抽取。

在知识图谱抽取过程的最后一步，这些实体和关系的实例成为图中独立的节点和边。每个节点和边的实体描述被聚合并摘要。关系被聚合为图的边，其中某个给定关系的重复次数成为边的权重。声明也以类似方式被聚合。在本文中，我们的分析对实体匹配使用精确字符串匹配——实体匹配指的是协调同一实体的不同抽取名称的任务（Barlaug and Gulla, 2021; Christen and Christen, 2012; Elmagarmid et al., 2006）。不过，只需对提示或代码稍作调整，也可以使用更宽松的匹配方法。此外，GraphRAG 对重复实体通常具有较强的鲁棒性，因为在后续步骤中，重复实体通常会被聚类在一起进行摘要。

\subsubsection{知识图谱 → 图社区}

给定上一步创建的图索引，可以使用多种社区检测算法将图划分为强连通节点构成的社区（例如，参见 Fortunato (2010) 和 Jin et al. (2021) 的综述）。在我们的流程中，我们以层级方式使用 Leiden 社区检测（Traag et al., 2019），在每个被检测到的社区内递归地检测子社区，直到到达无法再继续划分的叶子社区（leaf communities）。该层级结构的每一层都提供了一个社区划分，以互斥且共同穷尽（mutually exclusive, collectively exhaustive）的方式覆盖图的所有节点，从而支持分而治之（divide-and-conquer）的全局摘要。在某个示例数据集上进行此类层级划分的图示可参见附录~B。

\subsubsection{图社区 → 社区摘要}

下一步使用一种旨在能够扩展到超大规模数据集的方法，为社区层级结构中的每个社区创建类报告（report-like）的摘要。这些摘要本身就具有独立的用处，是理解数据集全局结构和语义的一种方式，甚至在没有特定查询的情况下也可用于对语料库进行意义建构。例如，用户可以浏览某一层级的社区摘要，寻找感兴趣的总体主题，然后阅读较低层级中链接的报告，以获得每个子主题的更多细节。不过，在此我们关注它们作为基于图的索引的一部分、用于回答全局查询时的用途。

GraphRAG 通过将各种元素摘要（针对节点、边及相关声明）添加到一个社区摘要模板中来生成社区摘要。来自较低层级社区的摘要被用于生成较高层级社区的摘要，具体方式如下：

\begin{itemize}[leftmargin=2em]
\item \textbf{叶子层级社区。} 叶子层级社区的元素摘要被排定优先级，然后迭代地添加到 LLM 上下文窗口中，直到达到 token 上限。优先级排定方式如下：对于每条社区边，按照其源节点与目标节点的合并度（即整体显著性）降序，依次添加源节点、目标节点、边本身以及相关声明的描述。
\item \textbf{较高层级社区。} 如果所有元素摘要都能放入上下文窗口的 token 上限内，则按叶子层级社区的方式处理，对社区内所有元素摘要进行摘要。否则，按子社区的元素摘要 token 数降序对子社区排序，并迭代地用（较短的）子社区摘要替换其对应的（较长的）元素摘要，直到能够放入上下文窗口内。
\end{itemize}

\subsubsection{社区摘要 → 社区答案 → 全局答案}

给定一个用户查询，可以在一个多阶段过程中使用上一步生成的社区摘要来生成最终答案。社区结构的层级性还意味着，可以使用来自不同层级的社区摘要来回答问题，从而引出一个问题：层级社区结构中某个特定层级是否能在摘要的细节程度与覆盖范围之间，为一般意义建构问题提供最佳平衡（在第~4~节中评估）。对于给定的社区层级，任意用户查询的全局答案按如下方式生成：

\begin{itemize}[leftmargin=2em]
\item \textbf{准备社区摘要。} 社区摘要被随机打乱，并被划分为预先指定 token 大小的若干块。这样可确保相关信息分散在各块之中，而不是集中（并可能丢失）于单个上下文窗口内。
\item \textbf{映射（map）社区答案。} 并行地生成中间答案。LLM 还被要求生成一个 0--100 之间的分数，表示所生成答案对回答目标问题的帮助程度。分数为 0 的答案会被过滤掉。
\item \textbf{归约（reduce）为全局答案。} 中间社区答案按帮助度分数降序排列，并迭代地添加到一个新的上下文窗口中，直到达到 token 上限。这个最终的上下文被用于生成返回给用户的全局答案。
\end{itemize}

\subsection{全局意义建构问题生成}

为了评估 RAG 系统在全局意义建构任务上的有效性，我们使用一个 LLM 生成一组语料库特定的问题，旨在评估对给定语料库的高层理解，而无需检索特定的低层事实。具体而言，给定对某个语料库及其用途的高层描述，LLM 被提示生成该 RAG 系统假想用户的人物画像。对于每个假想用户，LLM 随后被提示指明该用户会使用 RAG 系统完成的任务。最后，对于每个用户与任务的组合，LLM 被提示生成需要理解整个语料库的问题。算法~1 描述了该方法。

\begin{figure}[H]
\centering
\fbox{\begin{minipage}{0.93\linewidth}
\small
\textbf{算法 1：问题生成的提示流程}
\begin{enumerate}[leftmargin=2.2em,itemsep=2pt]
\item[1:] \textbf{输入}：某语料库的描述，用户数 $K$，每个用户的任务数 $N$，每个（用户，任务）组合的问题数 $M$。
\item[2:] \textbf{输出}：一组 $K \ast N \ast M$ 个需要对语料库进行全局理解的高层问题。
\item[3:] \textbf{procedure} \textsc{GenerateQuestions}
\item[4:] \quad 基于语料库描述，提示 LLM 执行以下操作：
\begin{enumerate}[leftmargin=1.5em,label=\arabic*.,itemsep=1pt]
\item 描述该数据集 $K$ 个潜在用户的人物画像。
\item 对每个用户，识别 $N$ 个与该用户相关的任务。
\item 针对每个用户与任务对，生成 $M$ 个满足以下条件的高层问题：
\begin{itemize}[leftmargin=1.2em,itemsep=0pt]
\item 需要对整个语料库的理解。
\item 不需要检索特定的低层事实。
\end{itemize}
\end{enumerate}
\item[5:] \quad 收集所生成的问题，为该数据集产生 $K \ast N \ast M$ 个测试问题。
\item[6:] \textbf{end procedure}
\end{enumerate}
\end{minipage}}
\end{figure}

在我们的评估中，我们设置 $K = M = N = 5$，每个数据集共计 125 个测试问题。表~1 展示了两个评估数据集各自的示例问题。

\subsection{评估全局意义建构的标准}

鉴于我们基于活动的意义建构问题缺乏金标准答案，我们采用一对一（head-to-head）比较方法，使用一个 LLM 评估者，依据特定标准来评判相对性能。我们设计了三个目标标准，用以刻画全局意义建构活动所期望的特质。附录~F 给出了我们使用 LLM 评估者计算一对一指标的提示，概括如下：

\begin{itemize}[leftmargin=2em]
\item \textbf{全面性（Comprehensiveness）。} 答案提供了多少细节来涵盖问题的所有方面与细节？
\item \textbf{多样性（Diversity）。} 答案在就问题提供不同视角和洞见方面有多丰富、多样？
\item \textbf{赋能性（Empowerment）。} 答案在帮助读者理解该主题并对其作出明智判断方面表现如何？
\end{itemize}

\begin{table}[H]\centering
\caption{LLM 基于目标数据集的简短描述所生成的潜在用户、任务和问题的示例。这些问题面向的是全局理解，而非特定细节。}
\label{tab:1}
\small
\renewcommand{\arraystretch}{1.25}
\begin{tabular}{@{}p{2.2cm}p{12.5cm}@{}}
\toprule
\textbf{数据集} & \textbf{示例活动框架与全局意义建构问题的生成} \\
\midrule
播客文字稿 & \textbf{用户}：一位寻求科技行业洞见与趋势的科技记者 \\
（Podcast transcripts） & \textbf{任务}：理解科技领袖如何看待政策与监管的作用 \\
 & \textbf{问题}： \\
 & 1. 哪些剧集主要涉及科技政策与政府监管？ \\
 & 2. 嘉宾如何看待隐私法对技术发展的影响？ \\
 & 3. 是否有嘉宾讨论创新与伦理考量之间的平衡？ \\
 & 4. 嘉宾提到了对当前政策的哪些建议性变更？ \\
 & 5. 是否讨论了科技公司与政府之间的合作，以及如何讨论的？ \\
\midrule
新闻文章 & \textbf{用户}：将时事融入课程的教育工作者 \\
（News articles） & \textbf{任务}：教授健康与养生相关内容 \\
 & \textbf{问题}： \\
 & 1. 哪些当前的健康主题可以被纳入健康教育课程？ \\
 & 2. 新闻文章如何论述预防医学与养生的概念？ \\
 & 3. 是否存在相互矛盾的健康类文章，如果有，原因是什么？ \\
 & 4. 基于新闻报道，可以对公共健康优先事项获得哪些洞见？ \\
 & 5. 教育工作者如何利用该数据集来强调健康素养的重要性？ \\
\bottomrule
\end{tabular}
\end{table}

此外，我们使用一个称为直接性（Directness）的"控制标准"（control criterion），它回答"答案在多大程度上具体而清晰地回应了问题？"。简单来说，直接性以一种适用于任何 LLM 生成式摘要的通用方式评估答案的简洁性。我们将其纳入，是为了使其充当一个参照，借此我们可以判断其他标准结果的合理性。由于直接性实际上与全面性和多样性相对立，我们不会期望任何方法能在全部四个标准上同时获胜。

在我们的评估中，LLM 会被提供问题以及来自两个竞争系统的生成答案，并被提示依据该标准比较两个答案，然后给出哪个答案更优的最终判定。LLM 要么指出一个获胜者；要么，如果两者本质上相似，则返回平局（tie）。为了考虑 LLM 生成固有的随机性，我们对每次比较运行多个重复实验，并在重复实验和问题之间对结果取平均。针对某个示例问题的答案，LLM 评估的图示可参见附录~D。

\section{分析}

\subsection{实验 1}

\subsubsection{数据集}

我们选取了两个百万 token 量级的数据集，每个数据集都代表了用户在真实世界活动中可能遇到的语料库：

\textbf{播客文字稿（Podcast transcripts）。} 《Behind the Tech with Kevin Scott》的公开文字稿，这是一档以微软首席技术官 Kevin Scott 与各科技领域思想领袖之间对话为内容的播客（Scott, 2024）。该语料库被划分为 1669 个 $\times$ 600-token 的文本块，块与块之间有 100-token 的重叠（约 100 万 token）。

\textbf{新闻文章（News articles）。} 一个由 2013 年 9 月至 2023 年 12 月期间发表的新闻文章组成的基准数据集，涵盖娱乐、商业、体育、科技、健康和科学等多个类别（Tang and Yang, 2024）。该语料库被划分为 3197 个 $\times$ 600-token 的文本块，块与块之间有 100-token 的重叠（约 170 万 token）。

\subsubsection{实验条件}

我们比较了六种条件，包括处于四个不同图社区层级（C0、C1、C2、C3）的 GraphRAG、一种将我们的 map-reduce 方法直接应用于源文本的文本摘要方法（TS），以及一种向量 RAG"语义检索"（semantic search）方法（SS）：

\begin{itemize}[leftmargin=2em]
\item \textbf{C0。} 使用根层级（数量最少）的社区摘要来回答用户查询。
\item \textbf{C1。} 使用高层级社区摘要来回答查询。如果存在，这些是 C0 的子社区；否则是向下投影的 C0 社区。
\item \textbf{C2。} 使用中间层级社区摘要来回答查询。如果存在，这些是 C1 的子社区；否则是向下投影的 C1 社区。
\item \textbf{C3。} 使用低层级（数量最多）的社区摘要来回答查询。如果存在，这些是 C2 的子社区；否则是向下投影的 C2 社区。
\item \textbf{TS。} 与第~3.1.6~节中相同的方法，区别在于 map-reduce 摘要阶段所打乱和分块的是源文本（而非社区摘要）。
\item \textbf{SS。} 向量 RAG 的一种实现，其中检索文本块并添加到可用的上下文窗口中，直到达到指定的 token 上限。
\end{itemize}

上下文窗口的大小以及用于生成答案的提示在所有六种条件下都相同（除了对引用样式作了细微修改，以匹配所使用上下文信息的类型）。各条件仅在上下文窗口内容的创建方式上有所不同。支持条件 C0--C3 的图索引是使用我们用于实体和关系抽取的通用提示创建的，其中实体类型和少样本示例针对数据所属领域进行了定制。

\subsubsection{配置}

我们为生成社区摘要、社区答案和全局答案使用了固定的 8k token 上下文窗口大小（在附录~C 中说明）。使用 600-token 窗口进行的图索引（在第~A.2~节中说明），对于播客数据集耗时 281 分钟，运行于一台虚拟机（16GB RAM，Intel(R) Xeon(R) Platinum 8171M CPU @ 2.60GHz），并使用面向 gpt-4-turbo 的公共 OpenAI 端点（2M TPM，10k RPM）。

我们使用 graspologic 库（Chung et al., 2019）实现了 Leiden 社区检测。用于生成图索引和全局答案的提示可参见附录~E，而用于根据我们的标准评估 LLM 响应的提示可参见附录~F。下一节所呈现结果的完整统计分析可参见附录~G。

\subsection{实验 2}

为了验证实验 1 得出的全面性和多样性结果，我们实现了这些特质的基于声明（claim-based）的度量。我们采用 Ni et al. (2024) 对事实声明的定义，即"显式呈现某些可验证事实的陈述"。例如，句子"加利福尼亚州和纽约州实施了可再生能源采用的激励措施，凸显了可持续性在政策决策中的广泛重要性"包含两个事实声明：（1）加利福尼亚州实施了可再生能源采用的激励措施；（2）纽约州实施了可再生能源采用的激励措施。

为了抽取事实声明，我们使用了 Claimify（Metropolitansky and Larson, 2025），这是一种基于 LLM 的方法，它识别答案中至少包含一个事实声明的句子，然后将这些句子分解为简单、自包含的事实声明。我们将 Claimify 应用于实验 1 各条件下生成的答案。在从每个答案中移除重复声明之后，我们抽取出 47,075 条唯一声明，平均每个答案 31 条声明。

我们定义了两个指标，数值越高表示性能越好：

\begin{enumerate}[leftmargin=2em]
\item \textbf{全面性}：度量为每种条件下生成的答案中所抽取声明的平均数量。
\item \textbf{多样性}：通过对每个答案的声明进行聚类并计算簇（cluster）的平均数量来度量。
\end{enumerate}

对于聚类，我们遵循 Padmakumar and He (2024) 所描述的方法，即使用 Scikit-learn 实现的凝聚式聚类（agglomerative clustering）（Pedregosa et al., 2011）。各簇通过"complete"链接方式进行合并，意味着只有当它们最远点之间的最大距离小于或等于预定义的距离阈值时，才会被合并。所使用的距离度量为 $1 - \text{ROUGE-L}$。由于距离阈值会影响簇的数量，我们报告了一系列阈值下的结果。

\section{结果}

\subsection{实验 1}

索引过程为播客数据集生成了一个由 8,564 个节点和 20,691 条边构成的图，为新闻数据集生成了一个更大的、由 15,754 个节点和 19,520 条边构成的图。表~2 展示了每个图社区层级结构在不同层级下的社区摘要数量。

\textbf{全局方法 vs.\ 向量 RAG。} 如图~2 和表~6 所示，全局方法在两个数据集的全面性和多样性两个标准上都显著优于传统向量 RAG（SS）。具体而言，对于播客文字稿，全局方法取得了 72--83\%（$p<.001$）的全面性胜率；对于新闻文章，取得了 72--80\%（$p<.001$）的全面性胜率；而多样性胜率则分别为 75--82\%（$p<.001$）和 62--71\%（$p<.01$）。我们将直接性用作有效性检验，证实了向量 RAG 在所有比较中产生了最直接的响应。

\textbf{赋能性。} 无论是全局方法对向量 RAG（SS），还是 GraphRAG 方法对源文本摘要（TS），赋能性比较都呈现出复杂混合的结果。使用 LLM 分析 LLM 针对该指标的推理表明，提供具体示例、引语和引用的能力被认为是帮助用户达成明智理解的关键。调优元素抽取提示或许有助于在 GraphRAG 索引中保留更多此类细节。

\begin{figure}[H]
\centering
\includegraphics[width=0.96\linewidth]{fig2.pdf}
\caption{在两个数据集、四个指标以及每次比较 125 个问题（每次重复五次并取平均）的条件下，（行条件）相对于（列条件）的一对一胜率百分比。每个数据集和指标的总体获胜者以粗体显示。自身胜率未计算，但显示为预期的 50\% 以供参考。所有 GraphRAG 条件在全面性和多样性上均优于朴素 RAG（na\"ive RAG）。在答案的全面性和多样性方面，条件 C1--C3 相比 TS（无图索引的全局文本摘要）也表现出轻微提升。}
\label{fig:2}
\end{figure}

\begin{table}[H]\centering
\caption{上下文单元（C0--C3 为社区摘要，TS 为文本块）的数量、对应的 token 计数，以及占最大 token 计数的百分比。对源文本进行 map-reduce 摘要是资源消耗最大的方法，需要最多的上下文 token。根层级社区摘要（C0）每次查询所需的 token 大幅减少（9$\times$--43$\times$）。}
\label{tab:2}
\small
\begin{tabular}{@{}l*{10}{c}@{}}
\toprule
& \multicolumn{5}{c}{\textbf{播客文字稿}} & \multicolumn{5}{c}{\textbf{新闻文章}} \\
\cmidrule(lr){2-6}\cmidrule(lr){7-11}
& C0 & C1 & C2 & C3 & TS & C0 & C1 & C2 & C3 & TS \\
\midrule
单元数 & 34 & 367 & 969 & 1310 & 1669 & 55 & 555 & 1797 & 2142 & 3197 \\
token 数 & 26657 & 225756 & 565720 & 746100 & 1014611 & 39770 & 352641 & 980898 & 1140266 & 1707694 \\
\% 最大值 & 2.6 & 22.2 & 55.8 & 73.5 & 100 & 2.3 & 20.7 & 57.4 & 66.8 & 100 \\
\bottomrule
\end{tabular}
\end{table}

\textbf{社区摘要 vs.\ 源文本。} 在使用 GraphRAG 将社区摘要与源文本进行比较时，除根层级摘要外，社区摘要总体上在答案的全面性和多样性方面带来了虽小但稳定的提升。播客数据集中的中间层级摘要和新闻数据集中的低层级社区摘要分别取得了 57\%（$p<.001$）和 64\%（$p<.001$）的全面性胜率。多样性胜率方面，播客中间层级摘要为 57\%（$p=.036$），新闻低层级社区摘要为 60\%（$p<.001$）。表~2 还说明了 GraphRAG 相比源文本摘要在可扩展性上的优势：对于低层级社区摘要（C3），GraphRAG 所需的上下文 token 减少了 26--33\%；而对于根层级社区摘要（C0），所需 token 减少了超过 97\%。与其他全局方法相比，根层级 GraphRAG 仅有适度的性能下降，却为意义建构活动所特有的迭代式问答提供了一种高效的方法，同时相比向量 RAG 仍保持了在全面性（72\% 胜率）和多样性（62\% 胜率）上的优势。

\begin{table}[H]\centering
\caption{按条件和数据集类型报告的所抽取声明的平均数量。加粗值表示每列中的最高分。}
\label{tab:3}
\small
\renewcommand{\arraystretch}{1.15}
\begin{tabular}{@{}lcc@{}}
\toprule
& \multicolumn{2}{c}{\textbf{声明的平均数量}} \\
\cmidrule(lr){2-3}
\textbf{条件} & \textbf{新闻文章} & \textbf{播客文字稿} \\
\midrule
C0 & \textbf{34.18} & 32.21 \\
C1 & 32.50 & 32.20 \\
C2 & 31.62 & \textbf{32.46} \\
C3 & 33.14 & 32.28 \\
TS & 32.89 & 31.39 \\
SS & 25.23 & 26.50 \\
\bottomrule
\end{tabular}
\end{table}

\subsection{实验 2}

表~3 展示了每种条件下所抽取声明平均数量（即基于声明的全面性度量）的结果。对于新闻和播客两个数据集，所有全局检索条件（C0--C3）和源文本摘要（TS）的全面性都高于向量 RAG（SS）。在所有情形下，这些差异都具有统计显著性（$p<.05$）。这些发现与实验 1 中基于 LLM 的胜率结果一致。

表~4 包含了簇平均数量（即基于声明的多样性度量）的结果。对于播客数据集，所有全局检索条件在所有距离阈值下的多样性都显著高于 SS（$p<.05$），这与实验 1 中观察到的胜率一致。然而，对于新闻数据集，只有 C0 在所有距离阈值下都显著优于 SS（$p<.05$）。尽管 C1--C3 的平均簇数也高于 SS，但这些差异仅在某些距离阈值下才具有统计显著性。在实验 1 中，所有全局检索条件在新闻数据集上都显著优于 SS——而不仅仅是 C0。不过，对于新闻数据集，SS 与全局检索条件之间平均多样性分数的差异要小于播客数据集，这在方向上与基于声明的结果相一致。

对于全面性和多样性两方面，在两个数据集上，全局检索条件之间或全局检索与 TS 之间均未观察到统计显著差异。

最后，对于实验 1 中的每一对成对比较，我们检验了 LLM 所偏好的答案是否与基于声明指标得出的获胜者一致。由于实验 1 中每对成对比较执行了五次，而基于声明的指标对每次比较只提供一个结果，我们使用多数投票（majority voting）将实验 1 的结果聚合为单一标签。例如，如果在某个问题的全面性上 C0 在五次判定中有三次胜过 SS，则 C0 被标记为获胜者，SS 为失败者。然而，如果 C0 胜两次、SS 胜一次、平局两次，则不存在多数结果，因此最终标签为平局。

我们发现，对于基于声明的指标，精确的平局十分罕见。一种可能的解决方案是基于阈值来定义平局（例如，条件 A 与条件 B 的基于声明的结果之间的绝对差必须小于或等于 $x$）。然而，我们观察到结果对阈值的选择很敏感。因此，我们聚焦于聚合后的 LLM 标签不是平局的情形，这些情形分别占全面性和多样性成对比较的 33\% 和 39\%。在这些情形下，聚合后的 LLM 标签与基于声明的标签在全面性上有 78\% 的成对比较相匹配，在多样性上有 69--70\%（跨所有距离阈值）相匹配，表明二者具有中等偏强的一致性。

\section{讨论}

\subsection{评估方法的局限性}

我们迄今的评估聚焦于针对两个语料库（每个约含 100 万 token）的意义建构问题。要理解性能如何泛化到来自不同领域、具有不同用例的数据集，还需要更多工作。此外，对杜撰率（fabrication rate）进行比较（例如使用 SelfCheckGPT（Manakul et al., 2023）之类的方法）也将加强当前的分析。

\begin{table}[H]\centering
\caption{按条件和数据集类型报告的不同距离阈值下簇的平均数量。加粗值表示每行中的最高分。}
\label{tab:4}
\small
\renewcommand{\arraystretch}{1.15}
\begin{tabular}{@{}llcccccc@{}}
\toprule
\textbf{数据集} & \textbf{距离阈值} & C0 & C1 & C2 & C3 & TS & SS \\
\midrule
\multirow{4}{*}{新闻文章}
 & 0.5 & \textbf{23.42} & 21.85 & 21.90 & 22.13 & 21.80 & 17.92 \\
 & 0.6 & \textbf{21.65} & 20.38 & 20.30 & 20.52 & 20.13 & 16.78 \\
 & 0.7 & \textbf{20.19} & 19.06 & 19.03 & 19.13 & 18.62 & 15.80 \\
 & 0.8 & \textbf{18.86} & 17.78 & 17.82 & 17.79 & 17.30 & 14.80 \\
\midrule
\multirow{4}{*}{播客文字稿}
 & 0.5 & \textbf{23.16} & 22.62 & 22.52 & 21.93 & 21.14 & 18.55 \\
 & 0.6 & \textbf{21.65} & 21.33 & 21.21 & 20.62 & 19.70 & 17.39 \\
 & 0.7 & \textbf{20.41} & 20.04 & 19.79 & 19.22 & 18.08 & 16.28 \\
 & 0.8 & \textbf{19.26} & 18.77 & 18.46 & 17.89 & 16.66 & 15.07 \\
\bottomrule
\end{tabular}
\end{table}

\subsection{未来工作}

支撑当前 GraphRAG 方法的图索引、丰富的文本注释和层级社区结构，为改进和适配提供了诸多可能性。这包括以更局部方式运作的 RAG 方法，即通过对用户查询与图注释进行基于嵌入的匹配。我们尤其看好混合式 RAG 方案，它将基于嵌入的匹配与即时（just-in-time）的社区报告生成相结合，然后再运用我们的 map-reduce 摘要机制。这种"上卷"（roll-up）方法还可以扩展到社区层级结构的多个层级，也可以实现为一种更具探索性的"下钻"（drill down）机制，沿着较高层级社区摘要中蕴含的信息线索（information scent）进行追踪。

\textbf{更广泛的影响。} 作为一种对大型文档集合进行问答的机制，如果所生成的答案不能准确表征源数据，那么下游的意义建构和决策任务就存在风险。系统的使用应当伴随对 AI 使用情况以及输出中潜在错误的明确披露。然而，与向量 RAG 相比，GraphRAG 在缓解这些下游风险方面展现出潜力，尤其是对于全局性质的问题——这类问题若用其他方式回答，可能会用检索到的事实样本冒充全局摘要而给出错误答案。

\section{结论}

我们提出了 GraphRAG，一种结合了知识图谱生成与查询聚焦摘要（QFS）的 RAG 方法，以支持人类对整个文本语料库进行意义建构。初步评估表明，相比向量 RAG 基线，它在答案的全面性和多样性两方面均带来显著提升，并且相比一种全局但无图（graph-free）的、使用 map-reduce 源文本摘要的方法也表现出有利的比较结果。对于需要在同一数据集上进行大量全局查询的情形，基于实体的图索引中根层级社区的摘要提供了一种数据索引，它既优于向量 RAG，又以一小部分的 token 成本取得了与其他全局方法相竞争的性能。

\section*{致谢}

我们还要感谢以下为本工作作出贡献的人员：Alonso Guevara Fernández、Amber Hoak、Andrés Morales Esquivel、Ben Cutler、Billie Rinaldi、Chris Sanchez、Chris Trevino、Christine Caggiano、David Tittsworth、Dayenne de Souza、Douglas Orbaker、Ed Clark、Gabriel Nieves-Ponce、Gaudy Blanco Meneses、Kate Lytvynets、Katy Smith、Mónica Carvajal、Nathan Evans、Richard Ortega、Rodrigo Racanicci、Sarah Smith 和 Shane Solomon。

\section*{参考文献}
\begin{enumerate}[leftmargin=2em]
\item Achiam, J., Adler, S., Agarwal, S., Ahmad, L., Akkaya, I., Aleman, F. L., Almeida, D., Altenschmidt, J., Altman, S., Anadkat, S., et al. (2023). Gpt-4 technical report. arXiv preprint arXiv:2303.08774.
\item Anil, R., Borgeaud, S., Wu, Y., Alayrac, J.-B., Yu, J., Soricut, R., Schalkwyk, J., Dai, A. M., Hauth, A., et al. (2023). Gemini: a family of highly capable multimodal models. arXiv preprint arXiv:2312.11805.
\item Baek, J., Aji, A. F., and Saffari, A. (2023). Knowledge-augmented language model prompting for zero-shot knowledge graph question answering. arXiv preprint arXiv:2306.04136.
\item Ban, T., Chen, L., Wang, X., and Chen, H. (2023). From query tools to causal architects: Harnessing large language models for advanced causal discovery from data.
\item Barlaug, N. and Gulla, J. A. (2021). Neural networks for entity matching: A survey. ACM Transactions on Knowledge Discovery from Data (TKDD), 15(3):1--37.
\item Baumel, T., Eyal, M., and Elhadad, M. (2018). Query focused abstractive summarization: Incorporating query relevance, multi-document coverage, and summary length constraints into seq2seq models. arXiv preprint arXiv:1801.07704.
\item Blondel, V. D., Guillaume, J.-L., Lambiotte, R., and Lefebvre, E. (2008). Fast unfolding of communities in large networks. Journal of statistical mechanics: theory and experiment, 2008(10):P10008.
\item Brown, T., Mann, B., Ryder, N., Subbiah, M., Kaplan, J. D., Dhariwal, P., Neelakantan, A., Shyam, P., Sastry, G., Askell, A., et al. (2020). Language models are few-shot learners. Advances in neural information processing systems, 33:1877--1901.
\item Cheng, X., Luo, D., Chen, X., Liu, L., Zhao, D., and Yan, R. (2024). Lift yourself up: Retrieval-augmented text generation with self-memory. Advances in Neural Information Processing Systems, 36.
\item Christen, P. and Christen, P. (2012). The data matching process. Springer.
\item Chung, J., Pedigo, B. D., Bridgeford, E. W., Varjavand, B. K., Helm, H. S., and Vogelstein, J. T. (2019). Graspy: Graph statistics in python. Journal of Machine Learning Research, 20(158):1--7.
\item Dang, H. T. (2006). Duc 2005: Evaluation of question-focused summarization systems. In Proceedings of the Workshop on Task-Focused Summarization and Question Answering, pages 48--55.
\item Elmagarmid, A. K., Ipeirotis, P. G., and Verykios, V. S. (2006). Duplicate record detection: A survey. IEEE Transactions on knowledge and data engineering, 19(1):1--16.
\item Es, S., James, J., Espinosa-Anke, L., and Schockaert, S. (2023). Ragas: Automated evaluation of retrieval augmented generation. arXiv preprint arXiv:2309.15217.
\item Etzioni, O., Cafarella, M., Downey, D., Kok, S., Popescu, A.-M., Shaked, T., Soderland, S., Weld, D. S., and Yates, A. (2004). Web-scale information extraction in knowitall: (preliminary results). In Proceedings of the 13th International Conference on World Wide Web, WWW '04, page 100--110, New York, NY, USA. Association for Computing Machinery.
\item Feng, Z., Feng, X., Zhao, D., Yang, M., and Qin, B. (2023). Retrieval-generation synergy augmented large language models. arXiv preprint arXiv:2310.05149.
\item Fortunato, S. (2010). Community detection in graphs. Physics reports, 486(3-5):75--174.
\item Gao, Y., Xiong, Y., Gao, X., Jia, K., Pan, J., Bi, Y., Dai, Y., Sun, J., and Wang, H. (2023). Retrieval-augmented generation for large language models: A survey. arXiv preprint arXiv:2312.10997.
\item He, X., Tian, Y., Sun, Y., Chawla, N. V., Laurent, T., LeCun, Y., Bresson, X., and Hooi, B. (2024). G-retriever: Retrieval-augmented generation for textual graph understanding and question answering. arXiv preprint arXiv:2402.07630.
\item Huang, J., Chen, X., Mishra, S., Zheng, H. S., Yu, A. W., Song, X., and Zhou, D. (2023). Large language models cannot self-correct reasoning yet. arXiv preprint arXiv:2310.01798.
\item Jacomy, M., Venturini, T., Heymann, S., and Bastian, M. (2014). Forceatlas2, a continuous graph layout algorithm for handy network visualization designed for the gephi software. PLoS ONE 9(6): e98679. \url{https://doi.org/10.1371/journal.pone.0098679}.
\item Jin, D., Yu, Z., Jiao, P., Pan, S., He, D., Wu, J., Philip, S. Y., and Zhang, W. (2021). A survey of community detection approaches: From statistical modeling to deep learning. IEEE Transactions on Knowledge and Data Engineering, 35(2):1149--1170.
\item Kang, M., Kwak, J. M., Baek, J., and Hwang, S. J. (2023). Knowledge graph-augmented language models for knowledge-grounded dialogue generation. arXiv preprint arXiv:2305.18846.
\item Khattab, O., Santhanam, K., Li, X. L., Hall, D., Liang, P., Potts, C., and Zaharia, M. (2022). Demonstrate-search-predict: Composing retrieval and language models for knowledge-intensive nlp. arXiv preprint arXiv:2212.14024.
\item Kim, D., Xie, L., and Ong, C. S. (2016). Probabilistic knowledge graph construction: Compositional and incremental approaches. In Proceedings of the 25th ACM International on Conference on Information and Knowledge Management, CIKM '16, page 2257--2262, New York, NY, USA. Association for Computing Machinery.
\item Kim, G., Kim, S., Jeon, B., Park, J., and Kang, J. (2023). Tree of clarifications: Answering ambiguous questions with retrieval-augmented large language models. arXiv preprint arXiv:2310.14696.
\item Klein, G., Moon, B., and Hoffman, R. R. (2006). Making sense of sensemaking 1: Alternative perspectives. IEEE intelligent systems, 21(4):70--73.
\item Kosinski, M. (2024). Evaluating large language models in theory of mind tasks. Proceedings of the National Academy of Sciences, 121(45):e2405460121.
\item Kuratov, Y., Bulatov, A., Anokhin, P., Sorokin, D., Sorokin, A., and Burtsev, M. (2024). In search of needles in a 11m haystack: Recurrent memory finds what llms miss.
\item LangChain (2024). Langchain graphs. \url{https://langchain-graphrag.readthedocs.io/en/latest/}.
\item Laskar, M. T. R., Hoque, E., and Huang, J. (2020). Query focused abstractive summarization via incorporating query relevance and transfer learning with transformer models. In Advances in Artificial Intelligence: 33rd Canadian Conference on Artificial Intelligence, Canadian AI 2020, Ottawa, ON, Canada, May 13--15, 2020, Proceedings 33, pages 342--348. Springer.
\item Lewis, P., Perez, E., Piktus, A., Petroni, F., Karpukhin, V., Goyal, N., K\"uttler, H., Lewis, M., Yih, W.-t., Rockt\"aschel, T., et al. (2020). Retrieval-augmented generation for knowledge-intensive nlp tasks. Advances in Neural Information Processing Systems, 33:9459--9474.
\item Liu, N. F., Lin, K., Hewitt, J., Paranjape, A., Bevilacqua, M., Petroni, F., and Liang, P. (2023). Lost in the middle: How language models use long contexts. arXiv:2307.03172.
\item LlamaIndex (2024). GraphRAG Implementation with LlamaIndex - V2. \url{https://github.com/run-llama/llama_index/blob/main/docs/docs/examples/cookbooks/GraphRAG_v2.ipynb}.
\item Madaan, A., Tandon, N., Gupta, P., Hallinan, S., Gao, L., Wiegreffe, S., Alon, U., Dziri, N., Prabhumoye, S., Yang, Y., et al. (2024). Self-refine: Iterative refinement with self-feedback. Advances in Neural Information Processing Systems, 36.
\item Manakul, P., Liusie, A., and Gales, M. J. (2023). Selfcheckgpt: Zero-resource black-box hallucination detection for generative large language models. arXiv preprint arXiv:2303.08896.
\item Mao, Y., He, P., Liu, X., Shen, Y., Gao, J., Han, J., and Chen, W. (2020). Generation-augmented retrieval for open-domain question answering. arXiv preprint arXiv:2009.08553.
\item Martin, S., Brown, W. M., Klavans, R., and Boyack, K. (2011). Openord: An open-source toolbox for large graph layout. SPIE Conference on Visualization and Data Analysis (VDA).
\item Melnyk, I., Dognin, P., and Das, P. (2022). Knowledge graph generation from text.
\item Metropolitansky, D. and Larson, J. (2025). Towards effective extraction and evaluation of factual claims.
\item Microsoft (2023). The impact of large language models on scientific discovery: a preliminary study using gpt-4.
\item Mooney, R. J. and Bunescu, R. (2005). Mining knowledge from text using information extraction. SIGKDD Explor. Newsl., 7(1):3--10.
\item NebulaGraph (2024). Nebulagraph launches industry-first graph rag: Retrieval-augmented generation with llm based on knowledge graphs. \url{https://www.nebula-graph.io/posts/graph-RAG}.
\item Neo4J (2024). Get started with graphrag: Neo4j's ecosystem tools. \url{https://neo4j.com/developer-blog/graphrag-ecosystem-tools/}.
\item Newman, M. E. (2006). Modularity and community structure in networks. Proceedings of the national academy of sciences, 103(23):8577--8582.
\item Ni, J., Shi, M., Stammbach, D., Sachan, M., Ash, E., and Leippold, M. (2024). AFaCTA: Assisting the annotation of factual claim detection with reliable LLM annotators. In Ku, L.-W., Martins, A., and Srikumar, V., editors, Proceedings of the 62nd Annual Meeting of the Association for Computational Linguistics (Volume 1: Long Papers), pages 1890--1912, Bangkok, Thailand. Association for Computational Linguistics.
\item OpenAI (2023). Chatgpt: Gpt-4 language model.
\item Padmakumar, V. and He, H. (2024). Does writing with language models reduce content diversity? ICLR.
\item Pedregosa, F., Varoquaux, G., Gramfort, A., Michel, V., Thirion, B., Grisel, O., Blondel, M., Prettenhofer, P., Weiss, R., Dubourg, V., Vanderplas, J., Passos, A., Cournapeau, D., Brucher, M., Perrot, M., and Duchesnay, E. (2011). Scikit-learn: Machine learning in python. Journal of Machine Learning Research, 12:2825--2830.
\item Ram, O., Levine, Y., Dalmedigos, I., Muhlgay, D., Shashua, A., Leyton-Brown, K., and Shoham, Y. (2023). In-context retrieval-augmented language models. Transactions of the Association for Computational Linguistics, 11:1316--1331.
\item Ranade, P. and Joshi, A. (2023). Fabula: Intelligence report generation using retrieval-augmented narrative construction. arXiv preprint arXiv:2310.13848.
\item Salminen, J., Liu, C., Pian, W., Chi, J., H\"ayh\"anen, E., and Jansen, B. J. (2024). Deus ex machina and personas from large language models: Investigating the composition of ai-generated persona descriptions. In Proceedings of the CHI Conference on Human Factors in Computing Systems, pages 1--20.
\item Sarthi, P., Abdullah, S., Tuli, A., Khanna, S., Goldie, A., and Manning, C. D. (2024). Raptor: Recursive abstractive processing for tree-organized retrieval. arXiv preprint arXiv:2401.18059.
\item Scott, K. (2024). Behind the Tech. \url{https://www.microsoft.com/en-us/behind-the-tech}.
\item Shao, Z., Gong, Y., Shen, Y., Huang, M., Duan, N., and Chen, W. (2023). Enhancing retrieval-augmented large language models with iterative retrieval-generation synergy. arXiv preprint arXiv:2305.15294.
\item Shin, J., Hedderich, M. A., Rey, B. J., Lucero, A., and Oulasvirta, A. (2024). Understanding human-ai workflows for generating personas. In Proceedings of the 2024 ACM Designing Interactive Systems Conference, pages 757--781.
\item Shinn, N., Cassano, F., Gopinath, A., Narasimhan, K., and Yao, S. (2024). Reflexion: Language agents with verbal reinforcement learning. Advances in Neural Information Processing Systems, 36.
\item Su, D., Xu, Y., Yu, T., Siddique, F. B., Barezi, E. J., and Fung, P. (2020). Caire-covid: A question answering and query-focused multi-document summarization system for covid-19 scholarly information management. arXiv preprint arXiv:2005.03975.
\item Tan, Z., Zhao, X., and Wang, W. (2017). Representation learning of large-scale knowledge graphs via entity feature combinations. In Proceedings of the 2017 ACM on Conference on Information and Knowledge Management, CIKM '17, page 1777--1786, New York, NY, USA. Association for Computing Machinery.
\item Tang, Y. and Yang, Y. (2024). MultiHop-RAG: Benchmarking retrieval-augmented generation for multi-hop queries. arXiv preprint arXiv:2401.15391.
\item Touvron, H., Martin, L., Stone, K., Albert, P., Almahairi, A., Babaei, Y., Bashlykov, N., Batra, S., Bhargava, P., Bhosale, S., et al. (2023). Llama 2: Open foundation and fine-tuned chat models. arXiv preprint arXiv:2307.09288.
\item Traag, V. A., Waltman, L., and Van Eck, N. J. (2019). From Louvain to Leiden: guaranteeing well-connected communities. Scientific Reports, 9(1).
\item Trajanoska, M., Stojanov, R., and Trajanov, D. (2023). Enhancing knowledge graph construction using large language models. ArXiv, abs/2305.04676.
\item Trivedi, H., Balasubramanian, N., Khot, T., and Sabharwal, A. (2022). Interleaving retrieval with chain-of-thought reasoning for knowledge-intensive multi-step questions. arXiv preprint arXiv:2212.10509.
\item Wang, J., Liang, Y., Meng, F., Sun, Z., Shi, H., Li, Z., Xu, J., Qu, J., and Zhou, J. (2023a). Is chatgpt a good nlg evaluator? a preliminary study. arXiv preprint arXiv:2303.04048.
\item Wang, S., Khramtsova, E., Zhuang, S., and Zuccon, G. (2024). Feb4rag: Evaluating federated search in the context of retrieval augmented generation. arXiv preprint arXiv:2402.11891.
\item Wang, X., Wei, J., Schuurmans, D., Le, Q., Chi, E., Narang, S., Chowdhery, A., and Zhou, D. (2022). Self-consistency improves chain of thought reasoning in language models. arXiv preprint arXiv:2203.11171.
\item Wang, Y., Lipka, N., Rossi, R. A., Siu, A., Zhang, R., and Derr, T. (2023b). Knowledge graph prompting for multi-document question answering.
\item Xu, Y. and Lapata, M. (2021). Text summarization with latent queries. arXiv preprint arXiv:2106.00104.
\item Yang, Z., Qi, P., Zhang, S., Bengio, Y., Cohen, W. W., Salakhutdinov, R., and Manning, C. D. (2018). HotpotQA: A dataset for diverse, explainable multi-hop question answering. In Conference on Empirical Methods in Natural Language Processing (EMNLP).
\item Yao, J.-g., Wan, X., and Xiao, J. (2017). Recent advances in document summarization. Knowledge and Information Systems, 53:297--336.
\item Yao, L., Peng, J., Mao, C., and Luo, Y. (2023). Exploring large language models for knowledge graph completion.
\item Yates, A., Banko, M., Broadhead, M., Cafarella, M., Etzioni, O., and Soderland, S. (2007). TextRunner: Open information extraction on the web. In Carpenter, B., Stent, A., and Williams, J. D., editors, Proceedings of Human Language Technologies: The Annual Conference of the North American Chapter of the Association for Computational Linguistics (NAACL-HLT), pages 25--26, Rochester, New York, USA. Association for Computational Linguistics.
\item Yuan, X., Li, J., Wang, D., Chen, Y., Mao, X., Huang, L., Xue, H., Wang, W., Ren, K., and Wang, J. (2024). S-eval: Automatic and adaptive test generation for benchmarking safety evaluation of large language models. arXiv preprint arXiv:2405.14191.
\item Zhang, J. (2023). Graph-toolformer: To empower llms with graph reasoning ability via prompt augmented by chatgpt. arXiv preprint arXiv:2304.11116.
\item Zhang, Y., Zhang, Y., Gan, Y., Yao, L., and Wang, C. (2024a). Causal graph discovery with retrieval-augmented generation based large language models. arXiv preprint arXiv:2402.15301.
\item Zhang, Z., Chen, J., and Yang, D. (2024b). Darg: Dynamic evaluation of large language models via adaptive reasoning graph. arXiv preprint arXiv:2406.17271.
\item Zheng, L., Chiang, W.-L., Sheng, Y., Zhuang, S., Wu, Z., Zhuang, Y., Lin, Z., Li, Z., Li, D., Xing, E., et al. (2024). Judging llm-as-a-judge with mt-bench and chatbot arena. Advances in Neural Information Processing Systems, 36.
\item Zhu, Y., Wang, X., Chen, J., Qiao, S., Ou, Y., Yao, Y., Deng, S., Chen, H., and Zhang, N. (2024). Llms for knowledge graph construction and reasoning: Recent capabilities and future opportunities.
\end{enumerate}

\appendix
\section{实体与关系抽取方法}

以下提示专为 GPT-4 设计，用于默认的 GraphRAG 初始化流程：

\begin{itemize}[leftmargin=2em]
\item 默认图抽取提示（Default Graph Extraction Prompt）
\item 声明抽取提示（Claim Extraction Prompt）
\end{itemize}

\subsection{实体抽取}

我们使用一个多部分 LLM 提示来完成此任务，它首先识别文本中的所有实体（包括其名称、类型和描述），然后识别明确相关实体之间的所有关系（包括源实体和目标实体以及对其关系的描述）。两类元素实例都以单个由分隔符分隔的元组列表形式输出。

\subsection{自我反思}

提示工程技术的选择对知识图谱抽取的质量有很强的影响（Zhu et al., 2024），并且不同的技术在模型所消耗和生成的 token 方面有不同的成本。自我反思（self-reflection）是一种提示工程技术，其中 LLM 先生成一个答案，然后被提示评估其输出的正确性、清晰度或完整性，最后基于该评估生成一个改进后的响应（Huang et al., 2023; Madaan et al., 2024; Shinn et al., 2024; Wang et al., 2022）。我们在知识图谱抽取中利用自我反思，并探索移除自我反思会如何影响性能和成本。

使用较大的文本块大小在 LLM 调用方面成本较低。然而，LLM 倾向于从较大尺寸的文本块中抽取较少的实体。例如，在一个示例数据集（HotPotQA，Yang et al., 2018）上，当文本块大小为 600 token 时，GPT-4 抽取的实体引用数量几乎是文本块大小为 2400 token 时的两倍。为了解决这一问题，我们部署了一种自我反思的提示工程方法。在从某个文本块抽取实体之后，我们将抽取出的实体反馈给 LLM，提示它"搜集"（glean）任何可能遗漏的实体。这是一个多阶段过程，其中我们首先要求 LLM 评估是否已抽取出所有实体，使用为 100 的 logit 偏置（logit bias）来强制其作出是/否决定。如果 LLM 回应说有实体被遗漏，那么一个指明"上次抽取中遗漏了许多实体"的续写会促使 LLM 检测这些被遗漏的实体。这种方法使我们能够在不降低质量（图~3）或不被迫引入噪声的情况下使用更大的文本块大小。我们将自我反思步骤迭代至指定的最大次数。

\begin{figure}[H]
\centering
\includegraphics[width=0.85\linewidth]{fig3.pdf}
\caption{对于我们使用 gpt-4-turbo 的通用实体抽取提示，在 HotPotQA 数据集（Yang et al., 2018）中检测到的实体引用数量如何随文本块大小和自我反思迭代次数而变化。图例为：600 文本块大小（600 chunk size）、1200 文本块大小（1200 chunk size）、2400 文本块大小（2400 chunk size）；纵轴为检测到的实体引用数（Entity references detected），横轴为执行的自我反思迭代次数（Number of self-reflection iterations performed）。}
\label{fig:3}
\end{figure}

\section{社区检测示例}

\begin{figure}[H]
\centering
\begin{subfigure}{0.48\linewidth}\centering
\includegraphics[width=\linewidth]{cand_p19_k00.pdf}
\caption{层级 0 的根社区（Root communities at level 0）}
\end{subfigure}\hfill
\begin{subfigure}{0.48\linewidth}\centering
\includegraphics[width=\linewidth]{cand_p19_k01.pdf}
\caption{层级 1 的子社区（Sub-communities at level 1）}
\end{subfigure}
\caption{在已索引的 MultiHop-RAG（Tang and Yang, 2024）数据集上，使用 Leiden 算法（Traag et al., 2019）检测到的图社区。圆圈表示实体节点，其大小与节点度成正比。节点布局通过 OpenORD（Martin et al., 2011）和 Force Atlas 2（Jacomy et al., 2014）完成。节点颜色表示实体社区，在两个层级的层级聚类下展示：(a) 层级 0，对应于具有最大模块度的层级划分；(b) 层级 1，揭示了这些根层级社区内部的结构。}
\label{fig:4}
\end{figure}

\section{上下文窗口选择}

上下文窗口大小对任何特定任务的影响尚不明确，尤其是对于像 gpt-4-turbo 这样具有 128k token 大型上下文规模的模型。鉴于信息有可能在较长上下文"中间丢失"（lost in the middle）（Kuratov et al., 2024; Liu et al., 2023），我们希望针对我们的数据集、问题和指标组合，探索改变上下文窗口大小所带来的影响。具体而言，我们的目标是确定基线条件（SS）的最优上下文大小，然后将其统一用于所有查询时的 LLM 使用。为此，我们测试了四种上下文窗口大小：8k、16k、32k 和 64k。令人意外的是，所测试的最小上下文窗口大小（8k）在全面性的所有比较上都普遍更优（平均胜率 58.1\%），同时在多样性（平均胜率 = 52.4\%）和赋能性（平均胜率 = 51.3\%）上与较大上下文大小表现相当。鉴于我们偏好更全面、更多样的答案，我们因此在最终评估中使用了固定的 8k token 上下文窗口大小。

\section{答案比较示例}

\begin{table}[H]\centering
\caption{针对新闻文章数据集的示例问题、答案以及 LLM 生成的评估。}
\label{tab:5}
\small
\renewcommand{\arraystretch}{1.2}
\begin{tabular}{@{}p{1.7cm}p{13cm}@{}}
\toprule
\textbf{问题} & 哪些公众人物在各类娱乐文章中被反复提及？ \\
\midrule
\textbf{Graph RAG} & 娱乐产业广阔而多元，涵盖电影、电视、音乐、体育和数字媒体。某些公众人物因其在这些领域的重大贡献和影响力而脱颖而出。以下摘要重点介绍了在各类娱乐文章中被反复提及的关键人物，反映了他们在该产业内的影响力与存在感。演员与导演 [...] 处于争议中的公众人物 [...] 音乐人与高管 [...] 运动员与教练 [...] 网红与创业者 [...] 这些人物在娱乐文章中被反复提及，标志着他们持续的相关性以及公众对其作品的兴趣。他们的影响力跨越娱乐的各个方面，从塑造影视中的文化叙事，到引领音乐和数字媒体的潮流。这些人物不仅为各自的领域作出贡献，还影响着更广泛的文化格局，常常成为社会讨论与公共话语的中心人物。 \\
\midrule
\textbf{朴素 RAG（Na\"ive RAG）} & 在各类娱乐文章中被反复提及的公众人物包括 Taylor Swift、Travis Kelce、Britney Spears 和 Justin Timberlake。这些人物因各种原因而受到关注，从他们的职业成就到他们的私人生活。Taylor Swift [...] Travis Kelce [...] Britney Spears [...] Justin Timberlake [...] 这些人物由于其高知名度以及公众对其职业生涯和私人生活的兴趣而被频繁报道。他们的活动，无论是在音乐、体育还是个人关系方面，都具有重大的文化和经济影响，正如媒体报道和公众反应所表明的那样。 \\
\midrule
\textbf{LLM 判定} & \textbf{全面性：获胜者=1（Graph RAG）。} 答案 1 更优，因为它提供了一份更全面、更详尽的公众人物清单，涵盖电影、电视、音乐、体育、游戏和数字媒体等更广泛的娱乐领域。它还包括他们贡献的示例以及他们对各自领域的影响，并提及了争议及其影响。答案 2 虽然在对少数几位人物的覆盖上较为详细，但仅限于较少数量的公众人物，且主要聚焦于他们的私人生活和关系，而非他们在整个娱乐产业中广泛的职业影响力。\newline
\textbf{多样性：获胜者=1（Graph RAG）。} 答案 1 更优，因为它通过涵盖来自娱乐产业不同领域（包括电影、电视、音乐、体育、游戏和数字媒体）的广泛公众人物，提供了更多样、更丰富的响应。它就这些人物的贡献和影响力，以及争议及其对公共话语的影响提供了洞见。该答案还为每位提及的人物引用了具体的数据来源，表明有多样的证据范围来支持其论断。相比之下，答案 2 聚焦于较小一组公众人物，主要来自音乐产业和体育界，并且严重依赖单一来源的数据，这使其在视角和洞见上不够多样。\newline
\textbf{赋能性：获胜者=1（Graph RAG）。} 答案 1 更优，因为它就娱乐产业各领域（包括电影、电视、音乐、体育和数字媒体）的公众人物提供了一份全面而有结构的概览。它列出了多位人物，提供了他们贡献的具体示例以及他们在娱乐文章中被提及的语境，并为每项论断附上了数据报告的引用。这种做法帮助读者理解该主题的广度，并在不被误导的情况下作出明智判断。相比之下，答案 2 聚焦于较小一组公众人物，且主要讨论他们的私人生活和关系，这可能无法对该主题提供同样广泛的理解。虽然答案 2 也引用了来源，但它在深度和多样性上不及答案 1。\newline
\textbf{直接性：获胜者=2（朴素 RAG）。} 答案 2 更优，因为它直接列出了在各类娱乐文章中被反复提及的具体公众人物，例如 Taylor Swift、Travis Kelce、Britney Spears 和 Justin Timberlake，并为他们被频繁提及的原因提供了简明的解释。答案 1 虽然全面，但包含了大量关于娱乐不同领域各类人物的详细信息，这些信息虽具信息量，却没有像答案 2 那样以同等程度的简洁性和具体性直接回答问题。 \\
\bottomrule
\end{tabular}
\end{table}

\section{系统提示}

\subsection{元素实例生成}

\begin{lstlisting}[style=promptstyle]
---Goal---
Given a text document that is potentially relevant to this activity and a list of entity types, identify all entities of those types from the text and all relationships among the identified entities.

---Steps---
1. Identify all entities. For each identified entity, extract the following information:
- entity_name: Name of the entity, capitalized
- entity_type: One of the following types: [{entity_types}]
- entity_description: Comprehensive description of the entity's attributes and activities

Format each entity as ("entity"{tuple_delimiter}<entity_name>{tuple_delimiter}<entity_type>{tuple_delimiter}<entity_description>

2. From the entities identified in step 1, identify all pairs of (source_entity, target_entity) that are *clearly related* to each other.
For each pair of related entities, extract the following information:
- source_entity: name of the source entity, as identified in step 1
- target_entity: name of the target entity, as identified in step 1
- relationship_description: explanation as to why you think the source entity and the target entity are related to each other
- relationship_strength: a numeric score indicating strength of the relationship between the source entity and target entity

Format each relationship as ("relationship"{tuple_delimiter}<source_entity>{tuple_delimiter}<target_entity>{tuple_delimiter}<relationship_description>{tuple_delimiter}<relationship_strength>)

3. Return output in English as a single list of all the entities and relationships identified in steps 1 and 2. Use **{record_delimiter}** as the list delimiter.

4. When finished, output {completion_delimiter}

---Examples---

Entity_types: ORGANIZATION,PERSON

Input:
The Fed is scheduled to meet on Tuesday and Wednesday, with the central bank planning to release its latest policy decision on Wednesday at 2:00 p.m. ET, followed by a press conference where Fed Chair Jerome Powell will take questions. Investors expect the Federal Open Market Committee to hold its benchmark interest rate steady in a range of 5.25%-5.5%.

Output:
("entity"{tuple_delimiter}FED{tuple_delimiter}ORGANIZATION{tuple_delimiter}The Fed is the Federal Reserve, which is setting interest rates on Tuesday and Wednesday)
{record_delimiter}
("entity"{tuple_delimiter}JEROME POWELL{tuple_delimiter}PERSON{tuple_delimiter}Jerome Powell is the chair of the Federal Reserve)
{record_delimiter}
("entity"{tuple_delimiter}FEDERAL OPEN MARKET COMMITTEE{tuple_delimiter}ORGANIZATION{tuple_delimiter}The Federal Reserve committee makes key decisions about interest rates and the growth of the United States money supply)
{record_delimiter}
("relationship"{tuple_delimiter}JEROME POWELL{tuple_delimiter}FED{tuple_delimiter}Jerome Powell is the Chair of the Federal Reserve and will answer questions at a press conference{tuple_delimiter}9)
{completion_delimiter}

...More examples...

---Real Data---

Entity_types: {entity_types}
Input: {input_text}
Output:
\end{lstlisting}

\subsection{社区摘要生成}

\begin{lstlisting}[style=promptstyle]
---Role---
You are an AI assistant that helps a human analyst to perform general information discovery. Information discovery is the process of identifying and assessing relevant information associated with certain entities (e.g., organizations and individuals) within a network.

---Goal---
Write a comprehensive report of a community, given a list of entities that belong to the community as well as their relationships and optional associated claims. The report will be used to inform decision-makers about information associated with the community and their potential impact. The content of this report includes an overview of the community's key entities, their legal compliance, technical capabilities, reputation, and noteworthy claims.

---Report Structure---
The report should include the following sections:
- TITLE: community's name that represents its key entities - title should be short but specific. When possible, include representative named entities in the title.
- SUMMARY: An executive summary of the community's overall structure, how its entities are related to each other, and significant information associated with its entities.
- IMPACT SEVERITY RATING: a float score between 0-10 that represents the severity of IMPACT posed by entities within the community. IMPACT is the scored importance of a community.
- RATING EXPLANATION: Give a single sentence explanation of the IMPACT severity rating.
- DETAILED FINDINGS: A list of 5-10 key insights about the community. Each insight should have a short summary followed by multiple paragraphs of explanatory text grounded according to the grounding rules below. Be comprehensive.

Return output as a well-formed JSON-formatted string with the following format:
{{
"title": <report_title>,
"summary": <executive_summary>,
"rating": <impact_severity_rating>,
"rating explanation": <rating_explanation>,
"findings": [
{{ "summary":<insight_1_summary>, "explanation": <insight_1_explanation> }},
{{ "summary":<insight_2_summary>, "explanation": <insight_2_explanation> }}
]
}}

---Grounding Rules---
Points supported by data should list their data references as follows:
"This is an example sentence supported by multiple data references [Data: <dataset name> (record ids); <dataset name> (record ids)]."

Do not list more than 5 record ids in a single reference. Instead, list the top 5 most relevant record ids and add "+more" to indicate that there are more.

For example:
"Person X is the owner of Company Y and subject to many allegations of wrongdoing [Data: Reports (1), Entities (5, 7); Relationships (23); Claims (7, 2, 34, 64, 46, +more)]."

where 1, 5, 7, 23, 2, 34, 46, and 64 represent the id (not the index) of the relevant data record.

Do not include information where the supporting evidence for it is not provided.

---Real Data---
Use the following text for your answer. Do not make anything up in your answer.
Input: {input_text}
...Report Structure and Grounding Rules Repeated...
Output:
\end{lstlisting}

\subsection{社区答案生成}

\begin{lstlisting}[style=promptstyle]
---Role---
You are a helpful assistant responding to questions about a dataset by synthesizing perspectives from multiple analysts.

---Goal---
Generate a response of the target length and format that responds to the user's question, summarize all the reports from multiple analysts who focused on different parts of the dataset, and incorporate any relevant general knowledge.

Note that the analysts' reports provided below are ranked in the **descending order of helpfulness**.

If you don't know the answer, just say so. Do not make anything up.

The final response should remove all irrelevant information from the analysts' reports and merge the cleaned information into a comprehensive answer that provides explanations of all the key points and implications appropriate for the response length and format.

Add sections and commentary to the response as appropriate for the length and format. Style the response in markdown.

The response shall preserve the original meaning and use of modal verbs such as "shall", "may" or "will".

The response should also preserve all the data references previously included in the analysts' reports, but do not mention the roles of multiple analysts in the analysis process.

Do not list more than 5 record ids in a single reference. Instead, list the top 5 most relevant record ids and add "+more" to indicate that there are more.

For example:
"Person X is the owner of Company Y and subject to many allegations of wrongdoing [Data: Reports (2, 7, 34, 46, 64, +more)]. He is also CEO of company X [Data: Reports (1, 3)]"

where 1, 2, 3, 7, 34, 46, and 64 represent the id (not the index) of the relevant data record.

Do not include information where the supporting evidence for it is not provided.

---Target response length and format---
{response_type}
---Analyst Reports---
{report_data}
...Goal and Target response length and format repeated...
Add sections and commentary to the response as appropriate for the length and format. Style the response in markdown.
Output:
\end{lstlisting}

\subsection{全局答案生成}

\begin{lstlisting}[style=promptstyle]
---Role---
You are a helpful assistant responding to questions about data in the tables provided.

---Goal---
Generate a response of the target length and format that responds to the user's question, summarize all relevant information in the input data tables appropriate for the response length and format, and incorporate any relevant general knowledge.

If you don't know the answer, just say so. Do not make anything up.

The response shall preserve the original meaning and use of modal verbs such as "shall", "may" or "will".

Points supported by data should list the relevant reports as references as follows:
"This is an example sentence supported by data references [Data: Reports (report ids)]"

Note: the prompts for SS (semantic search) and TS (text summarization) conditions use "Sources" in place of "Reports" above.

Do not list more than 5 record ids in a single reference. Instead, list the top 5 most relevant record ids and add "+more" to indicate that there are more.

For example:
"Person X is the owner of Company Y and subject to many allegations of wrongdoing [Data: Reports (2, 7, 64, 46, 34, +more)]. He is also CEO of company X [Data: Reports (1, 3)]"

where 1, 2, 3, 7, 34, 46, and 64 represent the id (not the index) of the relevant data report in the provided tables.

Do not include information where the supporting evidence for it is not provided.

At the beginning of your response, generate an integer score between 0-100 that indicates how **helpful** is this response in answering the user's question. Return the score in this format: <ANSWER_HELPFULNESS> score_value </ANSWER_HELPFULNESS>.

---Target response length and format---
{response_type}
---Data tables---
{context_data}
...Goal and Target response length and format repeated...
Output:
\end{lstlisting}

\section{评估提示}

\subsection{相对评估提示}

\begin{lstlisting}[style=promptstyle]
---Role---
You are a helpful assistant responsible for grading two answers to a question that are provided by two different people.

---Goal---
Given a question and two answers (Answer 1 and Answer 2), assess which answer is better according to the following measure:
{criteria}

Your assessment should include two parts:
- Winner: either 1 (if Answer 1 is better) and 2 (if Answer 2 is better) or 0 if they are fundamentally similar and the differences are immaterial.
- Reasoning: a short explanation of why you chose the winner with respect to the measure described above.

Format your response as a JSON object with the following structure:
{{
"winner": <1, 2, or 0>,
"reasoning": "Answer 1 is better because <your reasoning>."
}}

---Question---
{question}
---Answer 1---
{answer1}
---Answer 2---
{answer2}

Assess which answer is better according to the following measure:
{criteria}
Output:
\end{lstlisting}

\subsection{相对评估指标}

\begin{lstlisting}[style=promptstyle]
CRITERIA = {
"comprehensiveness": "How much detail does the answer provide to cover all the aspects and details of the question? A comprehensive answer should be thorough and complete, without being redundant or irrelevant. For example, if the question is 'What are the benefits and drawbacks of nuclear energy?', a comprehensive answer would provide both the positive and negative aspects of nuclear energy, such as its efficiency, environmental impact, safety, cost, etc. A comprehensive answer should not leave out any important points or provide irrelevant information. For example, an incomplete answer would only provide the benefits of nuclear energy without describing the drawbacks, or a redundant answer would repeat the same information multiple times.",

"diversity": "How varied and rich is the answer in providing different perspectives and insights on the question? A diverse answer should be multi-faceted and multi-dimensional, offering different viewpoints and angles on the question. For example, if the question is 'What are the causes and effects of climate change?', a diverse answer would provide different causes and effects of climate change, such as greenhouse gas emissions, deforestation, natural disasters, biodiversity loss, etc. A diverse answer should also provide different sources and evidence to support the answer. For example, a single-source answer would only cite one source or evidence, or a biased answer would only provide one perspective or opinion.",

"directness": "How specifically and clearly does the answer address the question? A direct answer should provide a clear and concise answer to the question. For example, if the question is 'What is the capital of France?', a direct answer would be 'Paris'. A direct answer should not provide any irrelevant or unnecessary information that does not answer the question. For example, an indirect answer would be 'The capital of France is located on the river Seine'.",

"empowerment": "How well does the answer help the reader understand and make informed judgements about the topic without being misled or making fallacious assumptions. Evaluate each answer on the quality of answer as it relates to clearly explaining and providing reasoning and sources behind the claims in the answer."
}
\end{lstlisting}

\section{统计分析}

表~\ref{tab:6} 给出了六种条件在四个指标上、跨 125 个问题和两个数据集的全部成对比较结果，详细的检验方法说明见表注。

\begin{center}
\scriptsize
\setlength{\tabcolsep}{4pt}
\renewcommand{\arraystretch}{1.1}
\begin{longtable}{@{}lll cccc cccc@{}}
\caption{在 125 个问题和两个数据集上，六种条件在四个指标上的成对比较。对于每个问题和指标，获胜条件得分为 100，失败条件得分为 0，平局时每个条件得分为 50。随后将这些分数在每个条件的五次评估运行上取平均。Shapiro--Wilk 检验结果表明数据不服从正态分布。因此，采用非参数检验（Wilcoxon 符号秩检验）来评估各对条件之间的性能差异，并应用 Holm--Bonferroni 校正以考虑多重成对比较。表明存在统计显著差异的校正后 $p$ 值以粗体突出显示。}
\label{tab:6}\\
\toprule
& & & \multicolumn{4}{c}{\textbf{播客文字稿}} & \multicolumn{4}{c}{\textbf{新闻文章}} \\
\cmidrule(lr){4-7}\cmidrule(lr){8-11}
\textbf{指标} & \textbf{条件 1} & \textbf{条件 2} & 均值 1 & 均值 2 & $Z$ 值 & $p$ 值 & 均值 1 & 均值 2 & $Z$ 值 & $p$ 值 \\
\midrule
\endfirsthead
\multicolumn{11}{c}{\textit{（表 6 续）}}\\
\toprule
& & & \multicolumn{4}{c}{\textbf{播客文字稿}} & \multicolumn{4}{c}{\textbf{新闻文章}} \\
\cmidrule(lr){4-7}\cmidrule(lr){8-11}
\textbf{指标} & \textbf{条件 1} & \textbf{条件 2} & 均值 1 & 均值 2 & $Z$ 值 & $p$ 值 & 均值 1 & 均值 2 & $Z$ 值 & $p$ 值 \\
\midrule
\endhead
\midrule
\multicolumn{11}{r}{\textit{（接下页）}}\\
\endfoot
\bottomrule
\endlastfoot
\multirow{15}{*}{全面性}
 & C0 & TS & 50.24 & 49.76 & $-0.06$ & 1 & 55.52 & 44.48 & $-2.03$ & 0.17 \\
 & C1 & TS & 51.92 & 48.08 & $-1.56$ & 0.633 & 58.8 & 41.2 & $-3.62$ & \textbf{0.002} \\
 & C2 & TS & 57.28 & 42.72 & $-4.1$ & \textbf{<0.001} & 62.08 & 37.92 & $-5.07$ & \textbf{<0.001} \\
 & C3 & TS & 56.48 & 43.52 & $-3.42$ & \textbf{0.006} & 63.6 & 36.4 & $-5.63$ & \textbf{<0.001} \\
 & C0 & SS & 71.92 & 28.08 & $-6.2$ & \textbf{<0.001} & 71.76 & 28.24 & $-6.3$ & \textbf{<0.001} \\
 & C1 & SS & 75.44 & 24.56 & $-7.45$ & \textbf{<0.001} & 74.72 & 25.28 & $-7.78$ & \textbf{<0.001} \\
 & C2 & SS & 77.76 & 22.24 & $-8.17$ & \textbf{<0.001} & 79.2 & 20.8 & $-8.34$ & \textbf{<0.001} \\
 & C3 & SS & 78.96 & 21.04 & $-8.12$ & \textbf{<0.001} & 79.44 & 20.56 & $-8.44$ & \textbf{<0.001} \\
 & TS & SS & 83.12 & 16.88 & $-8.85$ & \textbf{<0.001} & 79.6 & 20.4 & $-8.27$ & \textbf{<0.001} \\
 & C0 & C1 & 53.2 & 46.8 & $-1.96$ & 0.389 & 51.92 & 48.08 & $-0.45$ & 0.777 \\
 & C0 & C2 & 50.24 & 49.76 & $-0.23$ & 1 & 53.68 & 46.32 & $-1.54$ & 0.371 \\
 & C1 & C2 & 51.52 & 48.48 & $-1.62$ & 0.633 & 57.76 & 42.24 & $-4.01$ & \textbf{<0.001} \\
 & C0 & C3 & 49.12 & 50.88 & $-0.56$ & 1 & 52.16 & 47.84 & $-0.86$ & 0.777 \\
 & C1 & C3 & 50.32 & 49.68 & $-0.66$ & 1 & 55.12 & 44.88 & $-2.94$ & \textbf{0.016} \\
 & C2 & C3 & 52.24 & 47.76 & $-1.97$ & 0.389 & 58.64 & 41.36 & $-3.68$ & \textbf{0.002} \\
\midrule
\multirow{15}{*}{多样性}
 & C0 & TS & 50.24 & 49.76 & $-0.11$ & 1 & 46.88 & 53.12 & $-1.38$ & 0.676 \\
 & C1 & TS & 50.48 & 49.52 & $-0.12$ & 1 & 54.64 & 45.36 & $-1.88$ & 0.298 \\
 & C2 & TS & 57.12 & 42.88 & $-2.84$ & \textbf{0.036} & 55.76 & 44.24 & $-2.16$ & 0.184 \\
 & C3 & TS & 54.32 & 45.68 & $-2.39$ & 0.1 & 60.16 & 39.84 & $-4.07$ & \textbf{<0.001} \\
 & C0 & SS & 76.56 & 23.44 & $-7.12$ & \textbf{<0.001} & 62.08 & 37.92 & $-3.57$ & \textbf{0.003} \\
 & C1 & SS & 75.44 & 24.56 & $-7.33$ & \textbf{<0.001} & 64.96 & 35.04 & $-4.92$ & \textbf{<0.001} \\
 & C2 & SS & 80.56 & 19.44 & $-8.21$ & \textbf{<0.001} & 70.56 & 29.44 & $-6.29$ & \textbf{<0.001} \\
 & C3 & SS & 80.8 & 19.2 & $-8.3$ & \textbf{<0.001} & 69.12 & 30.88 & $-5.53$ & \textbf{<0.001} \\
 & TS & SS & 82.08 & 17.92 & $-8.43$ & \textbf{<0.001} & 67.2 & 32.8 & $-4.85$ & \textbf{<0.001} \\
 & C0 & C1 & 49.76 & 50.24 & $-0.13$ & 1 & 39.68 & 60.32 & $-3.61$ & \textbf{0.003} \\
 & C0 & C2 & 46.32 & 53.68 & $-1.5$ & 0.669 & 40.96 & 59.04 & $-3.14$ & \textbf{0.012} \\
 & C1 & C2 & 44.08 & 55.92 & $-3.27$ & \textbf{0.011} & 50.24 & 49.76 & $-0.22$ & 1 \\
 & C0 & C3 & 44 & 56 & $-2.6$ & 0.065 & 41.04 & 58.96 & $-3.47$ & \textbf{0.004} \\
 & C1 & C3 & 45.44 & 54.56 & $-2.98$ & \textbf{0.026} & 49.52 & 50.48 & $-0.01$ & 1 \\
 & C2 & C3 & 48.48 & 51.52 & $-0.96$ & 1 & 50.96 & 49.04 & $-0.39$ & 1 \\
\midrule
\multirow{15}{*}{赋能性}
 & C0 & TS & 40.96 & 59.04 & $-4.3$ & \textbf{<0.001} & 42.24 & 57.76 & $-3.32$ & \textbf{0.012} \\
 & C1 & TS & 45.2 & 54.8 & $-3.76$ & \textbf{0.002} & 50 & 50 & $-0.12$ & 1 \\
 & C2 & TS & 47.68 & 52.32 & $-2.2$ & 0.281 & 49.52 & 50.48 & $-0.22$ & 1 \\
 & C3 & TS & 48.72 & 51.28 & $-1.27$ & 1 & 51.68 & 48.32 & $-1.2$ & 1 \\
 & C0 & SS & 42.96 & 57.04 & $-3.71$ & \textbf{0.003} & 42.72 & 57.28 & $-3.12$ & \textbf{0.022} \\
 & C1 & SS & 47.68 & 52.32 & $-1.5$ & 0.936 & 51.36 & 48.64 & $-0.84$ & 1 \\
 & C2 & SS & 50.72 & 49.28 & $-0.55$ & 1 & 49.84 & 50.16 & $-0.2$ & 1 \\
 & C3 & SS & 48.96 & 51.04 & $-0.57$ & 1 & 49.52 & 50.48 & $-0.08$ & 1 \\
 & TS & SS & 57.52 & 42.48 & $-4.1$ & \textbf{<0.001} & 52.88 & 47.12 & $-1.1$ & 1 \\
 & C0 & C1 & 48.72 & 51.28 & $-1.23$ & 1 & 42.4 & 57.6 & $-3.9$ & \textbf{0.001} \\
 & C0 & C2 & 46.64 & 53.36 & $-2.54$ & 0.12 & 44.8 & 55.2 & $-2.16$ & 0.336 \\
 & C1 & C2 & 49.28 & 50.72 & $-1.73$ & 0.682 & 52 & 48 & $-1.45$ & 1 \\
 & C0 & C3 & 47.6 & 52.4 & $-1.78$ & 0.682 & 44.32 & 55.68 & $-3.45$ & \textbf{0.008} \\
 & C1 & C3 & 50 & 50 & 0 & 1 & 51.44 & 48.56 & $-1.02$ & 1 \\
 & C2 & C3 & 50.72 & 49.28 & $-0.86$ & 1 & 50.4 & 49.6 & $-0.22$ & 1 \\
\midrule
\multirow{15}{*}{直接性}
 & C0 & TS & 44.96 & 55.04 & $-4.09$ & \textbf{<0.001} & 45.2 & 54.8 & $-3.68$ & \textbf{0.003} \\
 & C1 & TS & 47.92 & 52.08 & $-2.41$ & 0.126 & 46.64 & 53.36 & $-2.91$ & \textbf{0.04} \\
 & C2 & TS & 48.8 & 51.2 & $-2.23$ & 0.179 & 48.32 & 51.68 & $-2.12$ & 0.179 \\
 & C3 & TS & 48.08 & 51.92 & $-2.23$ & 0.179 & 48.32 & 51.68 & $-2.56$ & 0.074 \\
 & C0 & SS & 35.12 & 64.88 & $-6.17$ & \textbf{<0.001} & 41.44 & 58.56 & $-4.82$ & \textbf{<0.001} \\
 & C1 & SS & 40.32 & 59.68 & $-4.83$ & \textbf{<0.001} & 45.2 & 54.8 & $-3.19$ & \textbf{0.017} \\
 & C2 & SS & 40.4 & 59.6 & $-4.67$ & \textbf{<0.001} & 44.88 & 55.12 & $-3.65$ & \textbf{0.003} \\
 & C3 & SS & 40.48 & 59.52 & $-4.69$ & \textbf{<0.001} & 45.6 & 54.4 & $-2.86$ & \textbf{0.043} \\
 & TS & SS & 43.6 & 56.4 & $-3.96$ & \textbf{<0.001} & 46 & 54 & $-2.68$ & 0.066 \\
 & C0 & C1 & 46.96 & 53.04 & $-2.87$ & \textbf{0.037} & 47.6 & 52.4 & $-2.17$ & 0.179 \\
 & C0 & C2 & 48.4 & 51.6 & $-2.06$ & 0.197 & 48.48 & 51.52 & $-1.61$ & 0.321 \\
 & C1 & C2 & 49.84 & 50.16 & $-1$ & 0.952 & 49.28 & 50.72 & $-1.6$ & 0.321 \\
 & C0 & C3 & 48.4 & 51.6 & $-1.8$ & 0.29 & 47.2 & 52.8 & $-2.62$ & 0.071 \\
 & C1 & C3 & 49.76 & 50.24 & 0 & 1 & 48.8 & 51.2 & $-1.29$ & 0.321 \\
 & C2 & C3 & 50 & 50 & 0 & 1 & 48.8 & 51.2 & $-1.84$ & 0.262 \\
\end{longtable}
\end{center}

\end{document}
